% 2020 北京卷 第 4 题：正三棱柱的三视图
% 来源 PDF：普通高考/2020/2020北京.pdf 第 1 页（A3 横排），印面 1113，第 4 题题干下方
% 原页矢量（pt，左上角为原点）：正（主）视图为 62.37 x 62.36 的正方形，中线一条竖直虚线；
%   侧（左）视图为同尺寸正方形；俯视图为底边 62.37、高 62.36 的等腰三角形（顶点在上）；
%   尺寸标注：左侧竖直双箭头「2」，主视图底边在中线处分为「1」「1」两段。
% 说明：原页三视图为示意画法（侧视图宽度与主视图相同、俯视图高等于底边），并非严格按
%   正三角形的高 sqrt3 缩放；题干另以「底面为正三角形」和尺寸 2、1、1 给出真实数据，
%   即底面正三角形边长 2、棱柱高 2。此处按原页如实重绘，不据图量取比例。
\begin{tikzpicture}[line width=0.7pt, font=\small,
  x=1.05cm, y=1.05cm, line cap=round, line join=round]

  % 正（主）视图：2 x 2 正方形 + 中线竖直虚线
  \draw (0,0) rectangle (2,2);
  \draw[dashed] (1,0) -- (1,2);

  % 侧（左）视图：2 x 2 正方形
  \draw (3,0) rectangle (5,2);

  % 俯视图：底边 2、顶点在上的等腰三角形（原页画法）
  \draw (0,-3.05) -- (2,-3.05) -- (1,-1.05) -- cycle;

  % 尺寸标注：高 2（左侧竖直双箭头）
  \draw[<->] (-0.30,0) -- (-0.30,2);
  \draw[thin] (-0.40,0) -- (0,0);
  \draw[thin] (-0.40,2) -- (0,2);
  \node[left, inner sep=3pt] at (-0.32,1) {$2$};

  % 尺寸标注：底边 1 与 1
  \draw[<->] (0,-0.24) -- (1,-0.24);
  \draw[<->] (1,-0.24) -- (2,-0.24);
  \node[below=1pt, inner sep=0pt] at (0.5,-0.24) {$1$};
  \node[below=1pt, inner sep=0pt] at (1.5,-0.24) {$1$};

  % 视图名称
  \node[below=2pt, inner sep=0pt] at (1,-0.72) {正（主）视图};
  \node[below=2pt, inner sep=0pt] at (4,-0.72) {侧（左）视图};
  \node[below=2pt, inner sep=0pt] at (1,-3.22) {俯视图};
\end{tikzpicture}
